\section{The compressed pencil}\label{sec:compressed}

Both observer conditions depend on the free rows only through the products $LN$ and $LAN$, where $N$ is a basis matrix of the common kernel of $C$ and $CA$. This section performs that compression and rewrites \eqref{eq:5a} and \eqref{eq:5b} in terms of the compressed data.

\subsection{The common kernel and the compression}\label{sec:compress}

Let
\begin{equation*}
K_0=\ker\begin{bmatrix}C\\ CA\end{bmatrix},\qquad k=\dim K_0,
\end{equation*}
and let $N\in\mathbb R^{n\times k}$ be a basis matrix for $K_0$. Define
\begin{equation*}
H=LN\in\mathbb R^{q\times k},\quad H_A=LAN\in\mathbb R^{q\times k},\quad \omega=\operatorname{rank}H .
\end{equation*}
The compressions below rest on a single rank identity.

\begin{lem}[Rank identity]\label{lem:rankid}
Let $P\in\mathbb C^{m\times n}$ and let $N_P\in\mathbb C^{n\times r}$ be a basis matrix for $\ker P$. Then, for every $M\in\mathbb C^{p\times n}$,
\begin{equation}\label{eq:rankid}
\operatorname{rank}\begin{bmatrix}P\\ M\end{bmatrix}
=\operatorname{rank}P+\operatorname{rank}(MN_P).
\end{equation}
\end{lem}

\begin{proof}
Set $\mathcal R_P=\operatorname{row}P$, $\mathcal R_M=\operatorname{row}M$ and $f(x)=xN_P$. Since $\operatorname{col}N_P=\ker P$,
\begin{equation*}
\ker f=(\ker P)^{\circ}=\mathcal R_P,\qquad
\operatorname{im}\bigl(f|_{\mathcal R_M}\bigr)=\operatorname{row}(MN_P),
\end{equation*}
where $(\cdot)^{\circ}$ denotes the left annihilator. Rank--nullity applied to $f|_{\mathcal R_M}$, whose kernel is $\mathcal R_M\cap\mathcal R_P$, gives
\begin{equation*}
\begin{aligned}
\dim(\mathcal R_P\cap\mathcal R_M)
&=\dim\ker\bigl(f|_{\mathcal R_M}\bigr)\\
&=\dim\mathcal R_M-\dim\operatorname{im}\bigl(f|_{\mathcal R_M}\bigr)\\
&=\operatorname{rank}M-\operatorname{rank}(MN_P),
\end{aligned}
\end{equation*}
and hence
\begin{equation*}
\begin{aligned}
\operatorname{rank}\begin{bmatrix}P\\ M\end{bmatrix}
&=\dim(\mathcal R_P+\mathcal R_M)\\
&=\operatorname{rank}P+\operatorname{rank}M-\dim(\mathcal R_P\cap\mathcal R_M)\\
&=\operatorname{rank}P+\operatorname{rank}(MN_P),
\end{aligned}
\end{equation*}
which is \eqref{eq:rankid}.
\end{proof}

\subsection{The compression and the pencil criterion}\label{sec:compression}

\begin{lem}[Compression and pencil criterion]\label{lem:compression}
With $H$, $H_A$ and $\omega$ as above, the following statements hold.

(i) Condition \eqref{eq:5a} holds if and only if $\operatorname{row}H_A\subseteq\operatorname{row}H$.

(ii) If \eqref{eq:5a} holds, then there exists $F\in\mathbb R^{q\times q}$ such that $H_A=FH$.

(iii) Assume \eqref{eq:5a}. Then \eqref{eq:5b} holds if and only if
\begin{equation*}
\operatorname{rank}(\lambda H-H_A)=\omega,\qquad \forall\lambda\in\mathbb C .
\end{equation*}
Equivalently, the $\omega\times\omega$ maximal minors of $P(\lambda)=\lambda H-H_A$ have no common complex root; equivalently, their greatest common divisor in $\mathbb C[\lambda]$ is a nonzero constant.

(iv) If \eqref{eq:5a} and \eqref{eq:5b} hold, then $\operatorname{rank}H_A=\operatorname{rank}H$.
\end{lem}

\begin{proof}
Set $P_0=[C;CA]$. Since $K_0=\ker P_0$, the matrix $N$ is a basis matrix for $\ker P_0$.

\emph{(i)} The rows of $P_0$ occur among the rows of each matrix below, so Lemma~\ref{lem:rankid} with $P=P_0$, applied once with $M=[LA;L]$ and once with $M=L$, gives
\begin{equation*}
\begin{aligned}
\operatorname{rank}\begin{bmatrix}LA\\ L\\ CA\\ C\end{bmatrix}
&=\operatorname{rank}P_0+\operatorname{rank}\begin{bmatrix}H_A\\ H\end{bmatrix},\\
\operatorname{rank}\begin{bmatrix}L\\ CA\\ C\end{bmatrix}
&=\operatorname{rank}P_0+\operatorname{rank}H .
\end{aligned}
\end{equation*}
The left-hand sides carry the same rows as the two matrices of \eqref{eq:5a}, so
\begin{equation*}
\eqref{eq:5a}
\iff
\operatorname{rank}\begin{bmatrix}H_A\\ H\end{bmatrix}=\operatorname{rank}H ,
\end{equation*}
which is equivalent to
\begin{equation}\label{eq:5a-1}
\operatorname{row}H_A\subseteq\operatorname{row}H .
\end{equation}

\emph{(ii)} By \eqref{eq:5a-1} every row of $H_A$ is a linear combination of the rows of $H$. Let $F$ have as its $i$-th row the coefficient vector of the $i$-th row of $H_A$; then
\begin{equation}\label{eq:5a-2}
H_A=FH .
\end{equation}

\emph{(iii)} Set
\begin{equation*}
M(\lambda)=[\lambda L-LA;CA;C],\qquad J=[L;CA;C].
\end{equation*}
Since $N$ annihilates both $C$ and $CA$,
\begin{equation*}
M(\lambda)N=\lambda H-H_A=P(\lambda),\qquad JN=H .
\end{equation*}
The rows of $P_0=[C;CA]$ occur among the rows of $M(\lambda)$ and of $J$, so that
\begin{equation*}
\operatorname{rank}\begin{bmatrix}P_0\\ M(\lambda)\end{bmatrix}=\operatorname{rank}M(\lambda),
\qquad
\operatorname{rank}\begin{bmatrix}P_0\\ J\end{bmatrix}=\operatorname{rank}J ,
\end{equation*}
and Lemma~\ref{lem:rankid} with $P=P_0$ yields
\begin{equation*}
\begin{aligned}
\operatorname{rank}M(\lambda)&=\operatorname{rank}P_0+\operatorname{rank}P(\lambda),\\
\operatorname{rank}J&=\operatorname{rank}P_0+\operatorname{rank}H .
\end{aligned}
\end{equation*}
Hence \eqref{eq:5b} is equivalent to
\begin{equation}\label{eq:5b-1}
\operatorname{rank}P(\lambda)=\operatorname{rank}H=\omega
\end{equation}
holding for every $\lambda\in\mathbb C$. Assume now \eqref{eq:5a}. By \eqref{eq:5a-2},
\begin{equation*}
P(\lambda)=\lambda H-H_A=(\lambda I-F)H ,
\end{equation*}
so every row of $P(\lambda)$ is a linear combination of the rows of $H$, that is, $\operatorname{row}P(\lambda)\subseteq\operatorname{row}H$, and therefore
\begin{equation}\label{eq:5b-2}
\operatorname{rank}P(\lambda)\le\omega,\qquad \forall\lambda\in\mathbb C .
\end{equation}
Combining \eqref{eq:5b-1} with the bound \eqref{eq:5b-2}, condition \eqref{eq:5b} becomes the requirement that some $\omega\times\omega$ minor of $P(\lambda)$ be nonzero at every $\lambda$:
\begin{equation*}
\eqref{eq:5b}
\iff
\forall\lambda\in\mathbb C\ \ \exists\,\omega\times\omega\text{ minor } f_i:\ f_i(\lambda)\neq0 .
\end{equation*}
Thus the $\omega\times\omega$ maximal minors $f_1,\dots,f_s$ of $P(\lambda)$ have no common complex root, equivalently their greatest common divisor in $\mathbb C[\lambda]$ is a nonzero constant. For $\omega=0$ the condition is vacuous.

\emph{(iv)} Assume \eqref{eq:5a} and \eqref{eq:5b}. By (iii), \eqref{eq:5b-1} holds for every $\lambda$; at $\lambda=0$, where $P(0)=-H_A$,
\begin{equation*}
\operatorname{rank}H_A=\operatorname{rank}P(0)=\operatorname{rank}H . \qedhere
\end{equation*}
\end{proof}

\section{The effective free-row dimension}\label{sec:dimension}

\subsection{The space of attainable row pairs}\label{sec:free}

Each free row $r_i\in\mathbb R^{1\times n}$ of $R$ determines the row pair
\begin{equation*}
(H_i,(H_A)_i)=r_i\,[N\ \ AN].
\end{equation*}
A free row is therefore specified by its coordinates relative to a basis of the subspace of attainable pairs
\begin{equation*}
\mathcal T=\operatorname{row}[N\ \ AN]\subseteq\mathbb R^{2k},\qquad
d_0=\dim\mathcal T=\operatorname{rank}[N\ \ AN],
\end{equation*}
and a rational basis of $\mathcal T$ represents each free row by exactly $d_0$ real coordinates.

\begin{lem}[Attainable row pairs and the effective dimension]\label{lem:attainable}
Every free row pair lies in $\mathcal T$, and writing the $r$ free rows in a basis of $\mathcal T$ represents them by
\begin{equation}\label{eq:kappa}
\kappa_{\rm obs}=r d_0
\end{equation}
coordinates. Moreover, when $r\ge1$,
\begin{equation}\label{lem3-2}
k\le d_0\le rd_0=\kappa_{\rm obs}.
\end{equation}
\end{lem}

\begin{proof}
The identity $(H_i,(H_A)_i)=r_i[N\ AN]$ places every free row pair in $\mathcal T$. Consider the linear map
\begin{equation*}
\Phi:\mathbb R^{1\times n}\to\mathcal T,\qquad \Phi(x)=x\,[N\ \ AN].
\end{equation*}
Its matrix relative to the standard basis of $\mathbb R^{1\times n}$ and a basis of $\mathcal T$ is $[N\ AN]$, so $\operatorname{rank}\Phi=\operatorname{rank}[N\ AN]$ and rank--nullity gives
\begin{equation*}
d_0=\dim\mathcal T=\dim\operatorname{im}\Phi=n-\dim\ker\Phi=\operatorname{rank}[N\ AN].
\end{equation*}
For the $r$ free rows set
\begin{equation*}
\Psi:(\mathbb R^{1\times n})^r\to\mathcal T^{\,r},\quad
\Psi(r_1,\dots,r_r)=\bigl(\Phi(r_1),\dots,\Phi(r_r)\bigr).
\end{equation*}
Its image is $\mathcal T^{\,r}$, so the number of coordinates after parametrization is
\begin{equation*}
\dim\mathcal T^{\,r}=r\dim\mathcal T=rd_0=\kappa_{\rm obs}.
\end{equation*}
Since $N$ has full column rank,
\begin{equation*}
\operatorname{rank}N=k,\qquad \operatorname{rank}[N\ AN]\ge\operatorname{rank}N=k,
\end{equation*}
that is,
\begin{equation}\label{lem3-1}
k\le d_0 .
\end{equation}
For $r\ge1$, the inequality $d_0\ge0$ gives
\begin{equation*}
d_0\le rd_0=\kappa_{\rm obs},
\end{equation*}
which together with \eqref{lem3-1} is \eqref{lem3-2}.
\end{proof}

\begin{remark}[Basis invariance]
If $N\mapsto NQ$ with $Q\in\mathbb R^{k\times k}$ nonsingular, then
\begin{equation*}
[N\ \ AN]\mapsto[N\ \ AN]\operatorname{diag}(Q,Q),
\end{equation*}
so $\mathcal T$ is unchanged and the coordinates of a free row transform by the nonsingular linear map $Q$. Hence $d_0=\dim\mathcal T$ and the number of parameters $\kappa_{\rm obs}=rd_0$ are basis-independent.
\end{remark}
